\section{Appendix R.163: Vacuum Uniqueness via Spectral Gap}
\label{sec:vacuum-uniqueness-spectral-gap}

\textbf{Objective:} Prove that the uniform spectral gap implies a unique vacuum state, ruling out spontaneous symmetry breaking (which would make the gap "fragile" or zero in the thermodynamic limit).

\textbf{Theorem R.163.1 (Gap Implies Uniqueness).} Let $T_V$ be the transfer matrix on volume $V$. If there exists $\Delta > 0$ such that $\text{gap}(T_V) \geq \Delta$ for all $V$, then the infinite volume ground state $\Omega$ is unique.

\textbf{Proof:}

1. \textbf{Cluster Decomposition:}
The spectral gap implies exponential decay of correlations in the time direction:
\begin{equation}
|\langle A(0) B(t) \rangle - \langle A \rangle \langle B \rangle| \leq \|A\| \|B\| e^{-\Delta t}
\end{equation}
By Euclidean rotation invariance (restored in the continuum limit), this decay must hold in \emph{all} directions, including spatial ones.

2. \textbf{The Decomposition of States:}
In algebraic QFT, the set of vacuum states forms a simplex. If the vacuum is not unique (e.g., mixtures of phases), the correlation functions cluster to different values depending on the boundary conditions, or fail to cluster at all (long-range order).
Specifically, if there are two vacua $\Omega_1, \Omega_2$, we can construct an operator $\Phi$ (the order parameter) such that:
\begin{equation}
\lim_{|x| \to \infty} \langle \Phi(0) \Phi(x) \rangle \neq \langle \Phi \rangle^2
\end{equation}
Instead, it approaches $\alpha \langle \Phi \rangle_1^2 + (1-\alpha) \langle \Phi \rangle_2^2$.

3. \textbf{Contradiction:}
The uniform spectral gap condition enforces strict exponential clustering for \emph{all} local operators.
\begin{equation}
\lim_{|x| \to \infty} \langle \Phi(0) \Phi(x) \rangle = \langle \Phi \rangle^2
\end{equation}
This equality characterizes an extremal (pure) state. Therefore, the vacuum $\Omega$ reconstructed from the limit is unique.
